%! The Normal Distribution and z-Scores — Unit 1, Lesson 1.7 — Student Packet Cover Sheet
% GRADUAL-RELEASE cover: four packet rows (Warm-Up · Guided Notes & Practice · AP Practice with
% NA · Homework). The EFFL spoiler rule is GONE — the learning targets name the formal
% vocabulary in bold, because the guided notes teach it directly today.
% The remindbox carries the lesson's IDEAS, never the lesson's process: no narrating the
% gradual release here (unit01/lesson03's cover is the model).
% \namedateperiod appears HERE ONLY — no name row on any other component.
\documentclass[12pt]{article}
\usepackage{apstats-article}
\usepackage{apstats-boxes}
\usepackage{ltablex}
\keepXColumns

\begin{document}

% Full-bleed navy banner. \coverbanner MEASURES the title block and sizes the
% band to it, so a lesson title long enough to wrap grows the band rather than
% running out the bottom of it.
\coverbanner{Unit 1}{Lesson 1.7 \quad The Normal Distribution and z-Scores}

\namedateperiod
\vspace{0.18in}

\boxguard[14]
\begin{learningtargetbox}
\small
% GRADUAL RELEASE: the vocabulary is taught in the guided notes, so the targets name it
% plainly and in bold. The AP Learning Objectives (VAR-2.A, VAR-2.B, VAR-2.C) stay in the
% teacher-facing plan.
\begin{enumerate}[leftmargin=1.5em, itemsep=5pt]
  \item \ldots{} tell a \textbf{parameter} ($\mu$, $\sigma$) from a \textbf{statistic}
        ($\bar x$, $s_x$), and say which one a summary is by asking who got measured.
  \item \ldots{} compute a \textbf{standardized score} (\textbf{$z$-score}) and interpret it in
        context --- how many standard deviations, on which side, of what.
  \item \ldots{} use the \textbf{empirical rule} on an approximately \textbf{normal}
        distribution to estimate what percent of values fall in a range, including a single tail.
  \item \ldots{} use $z$-scores to compare two results that came from two completely different
        distributions, and explain why the raw numbers cannot settle it.
\end{enumerate}
\end{learningtargetbox}

\vspace{0.14in}

\boxguard[14]
\begin{tocbox}
\small
\renewcommand{\arraystretch}{1.5}
\begin{tabularx}{\linewidth}{@{} c l X r @{}}
\rowcolor{sky}
\textbf{\#} & \textbf{Component} & \textbf{Description} & \textbf{Score} \\
\midrule
1 & Warm-Up & Two five-number summaries, counting steps away from an average, and one percent
                 split between two ends & \blank{1.2cm} \\
2 & Guided Notes \& Practice & Parameters and statistics, the $z$-score, the normal curve and
                 68--95--99.7 --- two tryouts, 500 loaves, 400 thermometers, and one loaf we
                 judge together & \blank{1.2cm} \\
% AP Practice is assigned but NOT scored: its score cell prints NA, never a \blank{}.
% Row 4 is the lesson's GRADED practice. Paper homework is the DEFAULT for this course; on the
% occasions DeltaMath is assigned instead, that replaces row 4 and its score cell is struck.
3 & AP Practice & Newborn birth weights: AP-style multiple choice and one free-response set ---
                 practice, not a grade & \textbf{NA} \\
4 & Homework & Parameter or statistic, three $z$-scores, two students on two different tests, a
                 normal curve to read, and one multiple-choice item ---
                 \textbf{scored, due the first class after two study halls} & \blank{1.2cm} \\
\midrule
  & \multicolumn{2}{r}{\textbf{Total}} & \blank{1.2cm} \\
\end{tabularx}
\end{tocbox}

\vspace{0.14in}

\boxguard[8]
\begin{remindbox}
\small
A raw number means nothing on its own. Twelve metres in the shot put and 12.3 seconds in the
100 m are not on the same scale, and no amount of staring at them will make them comparable ---
\textbf{so before you compare two results, count each one in the standard deviations of its own
distribution.} Keep the sign when you do: negative means below the mean, and below the mean is
good news in a race and bad news on a test. And the three percentages --- \textbf{68, 95,
99.7} --- are a fact about a \emph{shape}, not about data in general: they hold when a
distribution is mound-shaped and symmetric, and they hold for nothing else.
\end{remindbox}

\end{document}
